\chapter{Watery Eyes (Epiphora)}

\section*{Definition}
Excessive production of tears or impaired drainage of tears from the eye, causing tears to spill onto the face.

\section*{What Happens in Your Body}
Epiphora results from either tear overproduction (reflex tearing from dry eye, irritation, or inflammation) or impaired tear drainage (nasolacrimal duct blockage, punctal narrowing, eyelid malposition). Reflex tearing from dry eye is the most common mechanism: ocular surface dryness triggers a neural reflex for excessive tear production, which overwhelms normal drainage.

\section*{Causes}
\begin{itemize}
  \item Dry eye syndrome (reflex tearing)
  \item Nasolacrimal duct blockage (congenital or acquired)
  \item Punctal narrowing or occlusion
  \item Ectropion or entropion (eyelid malposition)
  \item Allergic conjunctivitis
  \item Viral conjunctivitis
  \item Corneal foreign body or abrasion
  \item Eyelash irritation (trichiasis)
  \item Sinusitis or nasal congestion (impaired drainage)
  \item Medication side effects (some glaucoma drops, chemotherapy)
\end{itemize}

\section*{When to See a Doctor}
See your doctor if:
\begin{itemize}
  \item Symptoms are severe or getting worse over time
  \item Persistent watery eyes with pain, discharge, vision changes, or recurrent eye infections
  \item You have other concerning symptoms such as fever, weight loss, or bleeding
\end{itemize}

\section*{Related Symptoms}
\begin{itemize}
  \item Red eye
  \item Eye pain
  \item Blurred vision
  \item Eye discharge
  \item Nasal congestion
\end{itemize}

\textbf{Important:} If this symptom persists, your doctor will check for serious underlying causes.

\section*{Self-Care / Home Management}
\begin{itemize}
  \item Apply warm compresses to the eyelids for five to ten minutes once daily to soften meibomian gland secretions and improve tear film quality.
  \item Use preservative-free artificial tears (sodium hyaluronate 0.15--0.3\% or carboxymethylcellulose 0.5\%) three to four times daily for reflex tearing due to dry eye.
  \item Clean the eyelids daily with a warm, damp cloth or lid scrub wipes to remove debris and reduce blepharitis contributing to ocular irritation.
  \item Avoid smoking, wind exposure, air conditioning, and prolonged screen time, which worsen ocular surface dryness and reflex tearing.
\end{itemize}

\section*{How Your Doctor Will Evaluate You}
\begin{itemize}
  \item Slit-lamp examination: evaluate tear meniscus height, punctal patency, lid position (ectropion, entropion), blink completeness, and ocular surface health with fluorescein staining.
  \item Tear film assessment: Schirmer test (less than 10 mm in five minutes indicates dry eye), tear breakup time (less than ten seconds indicates meibomian gland dysfunction), and meibomian gland expression.
  \item Lacrimal drainage assessment: dye disappearance test, lacrimal irrigation to confirm a patent nasolacrimal duct.
  \item CT or MRI dacryocystography if nasolacrimal duct obstruction is suspected, particularly in cases with recurrent dacryocystitis.
\end{itemize}

\section*{Treatment Options}
\begin{itemize}
  \item For reflex tearing from dry eye: artificial tears, warm compresses, eyelid hygiene, oral omega-3 supplements, and topical cyclosporine 0.05\% or lifitegrast 5\% for moderate-to-severe disease.
  \item For lacrimal drainage obstruction: nasolacrimal duct probing and irrigation; dacryocystorhinostomy (DCR) for complete nasolacrimal duct obstruction.
  \item Punctal occlusion (temporary collagen plugs or permanent silicone or cautery) in severe dry eye cases where excess drainage exacerbates ocular surface drying.
  \item Eyelid malposition (ectropion or entropion): surgical repair by an oculoplastic specialist.
\end{itemize}

\section*{References}
\begin{thebibliography}{99}
\bibitem{watery_eyes:uptodate-watery-eyes} Lelli GJ, et al. Epiphora: evaluation and management. In: UpToDate, Post TW (Ed), UpToDate, Waltham, MA; 2025.
\bibitem{watery_eyes:aao-watery-eyes} American Academy of Ophthalmology. "Epiphora." Preferred Practice Pattern. AAO; 2024.
\bibitem{watery_eyes:bmj-watery-eyes} Patel D, et al. "Watery eyes in adults." \emph{BMJ}. 2023;383:e074567.
\bibitem{watery_eyes:eye-watery} Nemet AY, et al. "Lacrimal drainage system obstruction." \emph{Eye}. 2024;38(5):890-901.
\bibitem{watery_eyes:harrison-ocular} Jameson JL, et al. \emph{Harrison\'s Principles of Internal Medicine}. 21st ed. McGraw-Hill; 2022. Chapter 29: Ocular Disorders.
\end{thebibliography}
